\chapter{Parameter estimation}
\label{ch:estimation}

With good data in hand, numbers have to be extracted from it. Two families of methods
exist, and the repository uses one now and is set up for the other.

\section{Linear regression form}

The surge equation \eqref{eq:surge} is nonlinear in $u$ but \emph{linear in its parameters}:
\begin{equation}
  \underbrace{T_k}_{y_k} =
  \underbrace{\begin{bmatrix} \dot u_k & u_k & |u_k|u_k & 1 \end{bmatrix}}_{\vect{\varphi}_k^\top}
  \underbrace{\begin{bmatrix} \mtot \\ X_u \\ \Xuu \\ c \end{bmatrix}}_{\vect{\theta}}
  + e_k ,
  \label{eq:regression}
\end{equation}
where $c$ is an optional constant offset (tether, current) and $e_k$ the equation error.
Stacking all samples gives $\vect{y} = \vect{\Phi}\vect{\theta} + \vect{e}$. This is exactly
what \code{surge_terms()} builds in \file{fitting/coupling_check.py}.

\section{Least squares}

The least-squares estimate minimises $\lVert\vect{y} - \vect{\Phi}\vect{\theta}\rVert^2$. Setting
the gradient to zero gives the \emph{normal equations}:
\begin{equation}
  \vect{\Phi}^\top\vect{\Phi}\,\hat{\vect{\theta}} = \vect{\Phi}^\top\vect{y}
  \quad\Rightarrow\quad
  \hat{\vect{\theta}} = (\vect{\Phi}^\top\vect{\Phi})^{-1}\vect{\Phi}^\top\vect{y}.
  \label{eq:ls}
\end{equation}
Geometrically, $\vect{\Phi}\hat{\vect{\theta}}$ is the orthogonal projection of $\vect{y}$
onto the column space of $\vect{\Phi}$, and the residual
$\vect{\varepsilon} = \vect{y} - \vect{\Phi}\hat{\vect{\theta}}$ is orthogonal to every
regressor. This fact is used again in \cref{ch:coupling}. In code,
\code{np.linalg.lstsq} solves \eqref{eq:ls} through an orthogonal factorisation, which is more
stable than forming $(\vect{\Phi}^\top\vect{\Phi})^{-1}$ explicitly.

\paragraph{Goodness of fit.} The coefficient of determination
$R^2 = 1 - \sum\varepsilon_k^2/\sum(y_k - \bar y)^2$ is the share of the variance in $\vect{y}$
explained by the model.

\paragraph{Uncertainty.} If the errors are white with variance $\sigma^2$, the covariance of the
estimate is
\begin{equation}
  \mathrm{Cov}(\hat{\vect{\theta}}) = \sigma^2(\vect{\Phi}^\top\vect{\Phi})^{-1}.
  \label{eq:cov}
\end{equation}
The diagonal gives the standard errors that \code{update_parameters()} can store in the
\code{uncertainty} field. A large entry flags a poorly excited parameter, because
$\vect{\Phi}^\top\vect{\Phi}$ is close to singular when two regressors move together.
\textbf{Caveat:} residuals of time series are rarely white. Neighbouring samples are correlated,
so the data contain fewer independent pieces of information than samples, and
\eqref{eq:cov} is \emph{too optimistic}. The effective-sample-size correction of
\cref{sec:neff} addresses the same problem for correlations.

\section{Equation error: the role of \texorpdfstring{$\dot u$}{du/dt}}

\Cref{eq:regression} needs the acceleration $\dot u$, which is not measured directly.
Differentiating a noisy signal amplifies noise in proportion to frequency (the
differentiator has gain $|j\omega| = \omega$). The code uses a \emph{Savitzky--Golay} (SG)
filter \cite{savitzky1964}: around each sample, a low-order polynomial is fitted by least
squares to a sliding window, and the smoothed value or its derivative is read off the
polynomial. Two properties make it suitable:
\begin{itemize}
  \item it is \textbf{zero-phase} (the window is centred), so it introduces no delay, which would
        otherwise bias the inertia as discussed in \cref{ch:timing};
  \item smoothing and differentiation come from the same polynomial fit, so $u$ and $\dot u$
        are mutually consistent.
\end{itemize}
\code{coupling_check.py} applies \emph{the same} SG filter (default window \qty{1}{\second},
quadratic) to every signal: thrust, velocities, attitude and PWM. If only some columns were
filtered, they would be smoothed differently, and the difference would appear as a spurious
term. Samples near stale-PWM gaps and near the ends of the record are dropped, because the
filter would mix bridged or missing data into them.

\subsection{How the Savitzky--Golay filter works}
\label{sec:sg-mechanics}

\begin{revised}[title=Added 16 October 2026,breakable]
\paragraph{Why the speed cannot be differentiated directly.} The EKF surge speed fluctuates from
sample to sample. In \file{zermatt_rectangle_08} the measured sample-to-sample standard deviation
is \qty{3.4}{\milli\metre\per\second}, at an EKF rate of about \qty{27}{\hertz}. (The first version
of this paragraph assumed \qty{0.03}{\metre\per\second}. That value was ten times too large and
was corrected on 21 October 2026, see \cref{sec:noise-budget}.) The finite difference
$(u_{k+1}-u_k)/\Delta t$ divides that fluctuation by a small $\Delta t$. At \qty{100}{\hertz}
it becomes an apparent acceleration of about \qty{0.3}{\metre\per\second\squared}, the same order
as the physical surge accelerations of \qtyrange{0.05}{0.3}{\metre\per\second\squared}. The
$\dot u$ column of $\vect{\Phi}$ would then be dominated by noise.

\paragraph{The procedure.} For every sample $k$, a window of $n$ neighbouring samples centred on
$k$ is taken, a polynomial of degree $p$ (here a parabola, $p = 2$) is fitted to them by least
squares, and the value of that polynomial at the centre is returned as the smoothed sample.
With the derivative option, the slope of the polynomial at the centre is returned instead. The
window then moves on by one sample. A parabola cannot follow fluctuations at the sample rate,
so these are averaged out, while slow changes pass through.

\paragraph{Equivalent weighted average.} Because the least-squares fit is linear in the data, the
result reduces to a fixed weighted sum of the window, with weights that depend only on $n$ and
$p$. For $n = 5$ and $p = 2$,
\begin{align}
  \hat u_k &= \tfrac{1}{35}\bigl(-3u_{k-2} + 12u_{k-1} + 17u_k + 12u_{k+1} - 3u_{k+2}\bigr),
  \label{eq:sg-smooth}\\
  \dot{\hat u}_k &= \tfrac{1}{10\,\Delta t}\bigl(-2u_{k-2} - u_{k-1} + u_{k+1} + 2u_{k+2}\bigr).
  \label{eq:sg-deriv}
\end{align}
The smoothing weights in \eqref{eq:sg-smooth} favour the centre and sum to one. The derivative
weights in \eqref{eq:sg-deriv} are antisymmetric, so they measure the tilt of the window and
vanish for a constant signal. \Cref{tab:sg-example} applies both to a short series. For the
centre sample, \eqref{eq:sg-smooth} gives
$(-3\cdot 2 + 12\cdot 2 + 17\cdot 5 + 12\cdot 4 - 3\cdot 5)/35 = 3.886$: the jump to 5 in the
raw data is pulled back towards its neighbours.

\paragraph{Why a polynomial and not a moving average.} A plain average over the window flattens
peaks and returns a poor slope. The local parabola keeps peaks and slopes approximately
correct, which matters because the slope $\dot u$ is the quantity of interest. The window is
centred on the sample, so no delay is introduced.

\paragraph{Ends of the record.} In the first and last half-window, the window extends past the
data. SciPy then fits the polynomial to the first or last full window instead, which makes these
values less reliable. They are removed by the \code{keep} mask.

\paragraph{Gaps.} The filter requires a continuous series without missing values. Thrust samples
with a stale PWM are therefore first replaced by linear interpolation between the valid samples
(\code{np.interp}), filtered together with the rest, and then excluded from the regression by the
\code{keep} mask, together with half a window on either side. The interpolated values never
enter the fit.
\end{revised}

\begin{incode}
\code{smooth(x, dt, window_s, poly=2, deriv=0)} in \file{fitting/coupling_check.py} converts the
window from seconds to samples, $n = \operatorname{round}(\code{window_s}/\Delta t)$, and forces it
to be odd (\code{| 1}), because the window needs a centre sample. At \qty{1}{\second} and
\qty{20}{\hertz}, $n = 21$. It then calls \code{scipy.signal.savgol_filter} with \code{delta=dt},
so that the derivative is returned in \si{\metre\per\second\squared} rather than per sample.
\end{incode}

\begin{table}[htbp]
  \centering
  \caption{Savitzky--Golay filter with a window of five samples and a quadratic polynomial,
    applied to a short series with unit sample spacing. Values computed with
    \code{scipy.signal.savgol_filter}.}
  \label{tab:sg-example}
  \begin{tabular}{lrrrrrrr}
    \toprule
    Sample $k$ & 0 & 1 & 2 & 3 & 4 & 5 & 6 \\
    \midrule
    Raw $u_k$ & 1 & 2 & 2 & 5 & 4 & 5 & 6 \\
    Smoothed $\hat u_k$ & 0.857 & 1.971 & 2.943 & 3.886 & 4.686 & 5.343 & 5.714 \\
    Slope $\dot{\hat u}_k$ & 1.186 & 1.043 & 0.900 & 0.800 & 0.800 & 0.514 & 0.229 \\
    \bottomrule
  \end{tabular}
\end{table}
\FloatBarrier

\begin{revised}[title=Added 14 October 2026: the window length biases the fit]
Zero phase is not the same as harmless. On a synthetic bag with perfect data, the
\qty{1}{\second} default window returns $\mtot$ \qty{14}{\percent} and $X_u$ \qty{18}{\percent}
too high, because the window is long against the surge transients ($\taueff \approx
\qty{0.2}{\second}$ with the dummy coefficients). A \qty{0.15}{\second} window removes that bias
on perfect data, but on noisy data it lets the noise of $\dot u$ through and underestimates
$\mtot$ (errors-in-variables, below). The window must be tuned against $\taueff$ and the noise
level, see \cref{tab:synth-window}. Whether the long default also affects the coupling
\emph{screen} (whose verdicts rest on correlations, not on the coefficients) has not been
checked yet. The synthetic bench cannot check it, because it simulates no coupling.
\end{revised}

\paragraph{Errors-in-variables.} Least squares assumes noise only in $\vect{y}$. Here the
regressors ($\dot u$ in particular) are also noisy, which biases the corresponding coefficient
toward zero (attenuation bias) \cite{greene2012}. Smoothing reduces this effect but cannot remove it entirely.

\section{Output error}
\label{sec:output-error}

The alternative is to simulate the model and compare the \emph{outputs}:
\begin{equation}
  \hat{\vect{\theta}} = \arg\min_{\vect{\theta}} \sum_k \bigl(u^\mathrm{meas}_k - u^\mathrm{sim}_k(\vect{\theta})\bigr)^2,
  \label{eq:output-error}
\end{equation}
where $u^\mathrm{sim}(\vect{\theta})$ comes from integrating \eqref{eq:surge} with the recorded
thrust. No derivative of noisy data is needed, a delay can be added as a parameter, and the
criterion is the one that matters for prediction. The cost is that \eqref{eq:output-error} is
nonlinear in $\vect{\theta}$ and needs an iterative solver (Gauss--Newton or
Levenberg--Marquardt, e.g.\ \code{scipy.optimize.least_squares}) and a reasonable starting
point. The equation-error estimate is the natural starting point \cite{klein2006}.

\begin{keyidea}
The two methods complement each other. Equation error is fast, linear and gives
diagnostics (residuals, partial correlations). Output error is robust to differentiation
noise and to delays, and it optimises the prediction directly. The pieces for output error
already exist: \code{rk4_step} is the simulator, and \code{propagate()} is almost
$u^\mathrm{sim}(\vect{\theta})$.
\end{keyidea}

\section{Verified on synthetic data}

\file{tests/test_coupling_check.py} generates \qty{300}{\second} of synthetic surge data with
known coefficients and white noise of \qty{0.3}{\newton}, and checks that the equation-error fit
recovers them. \Cref{tab:synthetic-fit} shows the result produced by the repository code.

\begin{table}[htbp]
\centering
\caption{Coefficients recovered by \code{own_axis_fit} from the synthetic data of
\file{tests/test_coupling_check.py}, without coupling ($k = 0$) and with an injected
coupling $k\,vr$, $k = -20$.}
\label{tab:synthetic-fit}
\begin{tabular}{@{}l S[table-format=3.1] S[table-format=3.2] S[table-format=3.2]@{}}
\toprule
Parameter & {True} & {Fit, $k=0$} & {Fit, $k=-20$} \\
\midrule
$\mtot$ / \unit{\kilogram} & 36.0 & 36.01 & 36.05 \\
$X_u$ / \unit{\newton\second\per\meter} & 8.0 & 7.92 & 7.95 \\
$\Xuu$ / \unit{\newton\second\squared\per\meter\squared} & 150.0 & 150.1 & 150.0 \\
offset / \unit{\newton} & 2.0 & 2.00 & 2.01 \\
$R^2$ & {--} & 0.99974 & 0.99943 \\
\bottomrule
\end{tabular}
\end{table}
\FloatBarrier

The coefficients are recovered even when the coupling term is present but missing from the
model. In this synthetic case $vr$ is uncorrelated with the own regressors, so leaving it out
does not bias them. It only lowers $R^2$ slightly. In real data the omitted term will rarely be
so cleanly uncorrelated, which is why it must be detected (\cref{ch:coupling}).

\section{Bias study of the surge fit}
\label{sec:bias-study}

\begin{revised}[title={New section, added 16 October 2026},breakable]
The first run of \file{fitting/fit_surge.py} on a perfect synthetic bag returned $\mtot$
\qty{8}{\percent} and $X_u$ \qty{13}{\percent} too high, and $\Xuu$ \qty{3}{\percent} too low,
at a \qty{20}{\hertz} grid and a \qty{1}{\second} smoothing window. That result raised three
questions: whether the smoothing alone causes such errors, whether the IMU accelerometer would be
the better source of $\dot u$, and how much a higher ArduSub publishing rate would help. This
section answers them. Every number is produced by \file{gen_biasstudy.py}. The script
\file{fit_surge.py} itself was not changed.
\end{revised}

Three variants of the equation-error fit of \eqref{eq:surge} are compared. Let $\mathcal{H}$
denote the Savitzky--Golay smoother with window $t_\mathrm{w}$.
\begin{description}[style=nextline]
  \item[A, current] As in \file{fit_surge.py}: $\mathcal{H}u$, the Savitzky--Golay derivative of
        $u$, $\mathcal{H}T$, and $|\mathcal{H}u|\,\mathcal{H}u$ as the quadratic regressor.
  \item[B, consistent] The same filter on every term, with the quadratic term formed from the raw
        speed before filtering and $\dot u$ taken as the derivative of the filtered speed:
        \begin{equation}
          \mathcal{H}T = \mtot\,\tfrac{\mathrm{d}}{\mathrm{d}t}(\mathcal{H}u)
            + X_u\,\mathcal{H}u + \Xuu\,\mathcal{H}(|u|u).
          \label{eq:consistent}
        \end{equation}
  \item[C, IMU] As B, but $\dot u$ replaced by $\mathcal{H}(f_x + g\sin\theta_0)$, the filtered
        accelerometer specific force with gravity removed.
\end{description}
The synthetic bags use the \code{steps} programme and the dummy T500 truth of
\cref{ch:synthetic}: \emph{ideal} (no imperfection), \emph{noise only} (default noise, no lags)
and \emph{all defaults} (noise, \qty{0.15}{\second} actuation lag, \qty{0.05}{\second} spin-up,
\qty{0.1}{\second} EKF lag).

\subsection{The idea in plain words}
\label{sec:plain-filter}

\begin{revised}[title={Added 16 October 2026},breakable]
\paragraph{Why smoothing is needed.} The fit balances, at every sample, thrust against mass
times acceleration plus drag. The acceleration is not recorded and has to be obtained from the
EKF speed. The speed wobbles by a few \unit{\centi\meter\per\second} between samples, and
dividing that wobble by \qty{0.05}{\second} produces a fake acceleration larger than the real one.
The speed is therefore smoothed first, by averaging each value with its neighbours over a window
of about a second. Savitzky--Golay is one recipe for that averaging.

\paragraph{The catch.} Smoothing spreads out every sudden change. When the thrust steps up, the
smoothed speed begins to rise slightly before the step and finishes slightly after it. If a
smoothed speed is compared with an unsmoothed thrust, the balance is wrong around every step,
and the least-squares fit distorts mass and drag to compensate.

\paragraph{The rule.} Every term is smoothed with the \emph{same} recipe. Smoothing is a
weighted average, and averages pass through sums and constant factors unchanged, so the balance
then still holds with the same coefficients. With mass 2, no drag and a three-point average:
\begin{center}
\small
\begin{tabular}{@{}lccccc@{}}
\toprule
Sample & 1 & 2 & 3 & 4 & 5 \\
\midrule
Acceleration & 0 & 0 & 3 & 0 & 0 \\
Thrust ($= 2\times$ acceleration) & 0 & 0 & 6 & 0 & 0 \\
Smoothed acceleration & 0 & 1 & 1 & 1 & 0 \\
Smoothed thrust & 0 & 2 & 2 & 2 & 0 \\
\bottomrule
\end{tabular}
\end{center}
The smoothed thrust is still exactly twice the smoothed acceleration. The fit returns the same
mass, without the noise.

\paragraph{Where the original script broke the rule.}
\begin{itemize}
  \item \emph{Acceleration.} The derivative mode of Savitzky--Golay returns the slope of each
        window's own parabola, which is a slightly different averaging from the one applied to the
        thrust. Now the speed is smoothed with the normal mode and the plain slope between
        neighbouring smoothed samples is taken (\code{np.gradient}). The slope of the smoothed
        speed equals the smoothed acceleration, so the recipe is the same.
  \item \emph{Quadratic drag.} Smoothing first and squaring afterwards is not the same as squaring
        first and smoothing afterwards. Speeds 0 and 2 give $1^2 = 1$ in the first case and
        $(0+4)/2 = 2$ in the second. The smoothed thrust contains the drag that actually acted,
        that is, an average of squares, so the drag column must be an average of squares too.
\end{itemize}
Savitzky--Golay itself was not replaced. It still performs the averaging. Only its separate
derivative mode is no longer used.
\end{revised}

\subsection{Two separate mechanisms on perfect data}

\paragraph{Filter inconsistency.} A linear time-invariant filter commutes with every linear
operation. If all four terms of the surge equation pass through the same $\mathcal{H}$, the
filtered signals satisfy \eqref{eq:consistent} \emph{exactly}, whatever the window, and the
filter introduces no bias on noise-free data. Variant A breaks this in two places. First, the
Savitzky--Golay derivative filter of \eqref{eq:sg-deriv} is not the derivative of the smoothing
filter \eqref{eq:sg-smooth}. Their frequency responses differ by more than the factor $j\omega$,
so $\dot u$ and $T$ are effectively filtered differently. Second, $|\mathcal{H}u|\,\mathcal{H}u
\neq \mathcal{H}(|u|u)$, because the quadratic term is nonlinear. Both errors grow with the
window and are largest at the steps, which is exactly where $\mtot$ is determined.
\Cref{tab:bias-filter} separates them from the grid effect by using a \qty{200}{\hertz} grid.

\begin{table}[htbp]
\centering
\caption{Relative error in percent of the surge coefficients fitted to the ideal synthetic bag
on a \qty{200}{\hertz} grid, for variant A (current) and variant B (consistent filtering), against
the smoothing window $t_\mathrm{w}$.}
\label{tab:bias-filter}
\small
\begin{tabular}{@{}c ccc ccc@{}}
\toprule
& \multicolumn{3}{c}{A, current} & \multicolumn{3}{c}{B, consistent} \\
\cmidrule(lr){2-4}\cmidrule(l){5-7}
$t_\mathrm{w}$ (\unit{\second}) & $\mtot$ & $X_u$ & $\Xuu$ & $\mtot$ & $X_u$ & $\Xuu$ \\
\midrule
\inputrows data/bias_filter_rows.tex
\bottomrule
\end{tabular}
\end{table}
\FloatBarrier

\paragraph{Grid quantisation of the thrust step.} The thrust is reconstructed from a held PWM,
so it jumps at the logged edge. On the analysis grid the jump appears at the first grid point
after the edge, between $0$ and $1/f_\mathrm{grid}$ late, on average $1/(2f_\mathrm{grid})$.
The speed, being continuous, is interpolated without that delay. The thrust column therefore
lags the speed by \qty{25}{\milli\second} on average at \qty{20}{\hertz}. That is a small fraction
of $\taueff \approx \qty{0.2}{\second}$, but it acts like the timing errors of \cref{ch:timing}
and lowers $\mtot$, as \cref{tab:bias-grid} shows.

\begin{table}[htbp]
\centering
\caption{Relative error in percent of $\mtot$ and $X_u$ fitted to the ideal synthetic bag with a
\qty{1}{\second} window, against the analysis grid rate $f_\mathrm{grid}$, for the three variants
(only $\mtot$ for C).}
\label{tab:bias-grid}
\small
\begin{tabular}{@{}c cc cc c@{}}
\toprule
& \multicolumn{2}{c}{A, current} & \multicolumn{2}{c}{B, consistent} & C, IMU \\
\cmidrule(lr){2-3}\cmidrule(lr){4-5}\cmidrule(l){6-6}
$f_\mathrm{grid}$ (\unit{\hertz}) & $\mtot$ & $X_u$ & $\mtot$ & $X_u$ & $\mtot$ \\
\midrule
\inputrows data/bias_grid_rows.tex
\bottomrule
\end{tabular}
\end{table}
\FloatBarrier

\begin{keyidea}
The \qty{8}{\percent} error of $\mtot$ at \qty{20}{\hertz} and \qty{1}{\second} is the
\emph{sum} of two biases of opposite sign: the filter inconsistency (about
\qty{+18}{\percent}) and the grid quantisation (about \qty{-8}{\percent}). Smoothing alone is
worse than the first result suggested. Both biases can be removed without new data: consistent
filtering \eqref{eq:consistent} removes the first at any window, and a grid of
\qtyrange{100}{200}{\hertz} reduces the second to one or two percent. The window then only trades
noise against resolution and no longer biases the fit on clean data.
\end{keyidea}

\subsection{Is the IMU a better source of \texorpdfstring{$\dot u$}{du/dt}?}

\paragraph{No acceleration from the EKF (added 16 October 2026).} The EKF does not estimate
acceleration. \code{publish_acceleration} is \code{false} in both \file{ekf_local.yaml} and
\file{ekf_global.yaml}. The global configuration also locks the acceleration states at zero
with a tiny initial covariance, because nothing measures them and they drifted otherwise. Even
if published, an EKF acceleration would be a filtered difference of velocities, with the same
problem as above plus an unknown lag. The only direct measurement of acceleration is the IMU.
The bags also contain \code{/filter/free_acceleration}, the Xsens' own gravity-free
acceleration, computed with its internal attitude estimate. It may avoid the gravity-leakage
problem discussed below. It is expressed in the world frame, so it must be rotated by the heading
to give the surge component, and its timestamps have the same question mark. It has not been
tested yet.

\paragraph{Noise budget (added 21 October 2026).}
\label{sec:noise-budget}
The sensor noise has been measured in this repository.
\file{src/sensors/dvl_a50_pkg/dvl_a50_pkg/measurement_noise_constants.py} and the docstrings of
\file{scripts/stationary_noise_comparison.py} and \file{scripts/stationary_allan_variance.py}
give values from a 64-minute stationary recording (2026-04-19). The POLARIS final report
(\file{scripts/sysid/Final_Report_POLARIS.pdf}, Section 5.3.1 and Appendix E.7) describes the setup:
the DVL was removed from the vehicle and fixed rigidly in the water, and the IMU was logged on the
vehicle resting on land. The report places the recording at Lake Zurich, while the script docstrings
say St.~Moritz. The values are:
\begin{itemize}
  \item IMU (Xsens Sirius, about \qty{86}{\hertz}): accelerometer $x$ standard deviation
        \qty{0.0144}{\metre\per\second\squared};
  \item DVL $v_x$ standard deviation \qtyrange{2.1}{7.5}{\milli\metre\per\second}
        (\code{stationary_11} and \code{stationary_02}).
\end{itemize}
On the vehicle, the EKF surge speed of \file{zermatt_rectangle_08} has a noise of about
\qty{1.5}{\milli\metre\per\second} around a one-second running median. The EKF speed is driven
by the DVL at about \qty{8}{\hertz}, so a window $t_\mathrm{w}$ holds roughly $8\,t_\mathrm{w}$
independent samples. The least-squares slope over $n$ independent samples spaced across
$t_\mathrm{w}$ has a standard deviation of about $\sigma_u\sqrt{12/n}/t_\mathrm{w}$. The averaged
accelerometer has about $\sigma_a/\sqrt{86\,t_\mathrm{w}}$. \Cref{tab:noise-budget} compares the
two.

\begin{table}[htbp]
\centering
\caption{Estimated random error of $\dot u$ for the differentiated EKF speed and for the averaged
accelerometer against the window $t_\mathrm{w}$, from the measured noise levels.}
\label{tab:noise-budget}
\small
\begin{tabular}{@{}ccc@{}}
\toprule
$t_\mathrm{w}$ (\unit{\second}) & EKF speed, differentiated (\unit{\milli\metre\per\second\squared}) & IMU, averaged (\unit{\milli\metre\per\second\squared}) \\
\midrule
0.25 & 14.7 & 3.1 \\
0.5  & 5.2  & 2.2 \\
1.0  & 1.8  & 1.6 \\
2.0  & 0.6  & 1.1 \\
\bottomrule
\end{tabular}
\end{table}
\FloatBarrier

For windows of half a second or more, both routes give random errors of a few
\unit{\milli\metre\per\second\squared}. Against step accelerations of
\qtyrange{50}{300}{\milli\metre\per\second\squared}, noise is then not what decides between them. The
final report, however, measured much smaller accelerations in closed loop. Its Table 6.2 gives a
maximum trackable surge acceleration of only \qty{0.02}{\metre\per\second\squared} from ramp
responses in Lake Zurich. That is the acceleration the velocity controller achieved, not the
open-loop step response, which should be larger. If the real step accelerations are nearer
\qty{20}{\milli\metre\per\second\squared}, the \qty{5}{\milli\metre\per\second\squared} of a
\qty{0.5}{\second} window is a quarter of the signal and biases $\mtot$ low through errors in
variables. A window of \qty{1}{\second} or more is then needed for the speed route. The first clean
step bag will settle this. The IMU wins clearly only below about \qty{0.3}{\second}, where the
differentiated speed grows as $t_\mathrm{w}^{-3/2}$. Short windows matter only if $\taueff$ is
short. The systematic errors are what decide. A pitch error of \qty{1}{\degree} leaks
\qty{171}{\milli\metre\per\second\squared} of gravity into the accelerometer, a hundred times its
averaged noise, and the IMU timestamps arrive in bursts. With the measured noise levels, the
Savitzky--Golay route on the EKF speed is therefore the better choice for now. Three caveats
apply. The stationary IMU figure was measured on the vehicle at rest on land, with the thrusters
off, not under thruster vibration (\qty{0.42}{\metre\per\second\squared} above \qty{5}{\hertz} on the vehicle, which the
window removes). The Allan-variance output (\file{stationary11_allan_params.json}) is not in the
repository. The low EKF speed noise is partly bought with the EKF's own lag.

\paragraph{Measurement, not integration.} Integrating the accelerometer to obtain a speed does
not help. A bias of \qty{0.01}{\meter\per\second\squared} alone gives a drift of
\qty{0.6}{\meter\per\second} after one minute, and fusing the IMU with the DVL and GNSS is what
the EKF already does. The useful idea is the opposite one: the accelerometer \emph{measures}
$\dot u$ directly, so nothing has to be differentiated. In FLU, with pitch positive nose-down
and neglecting lever-arm terms, $\dot u = f_x + g\sin\theta$ (\cref{ch:synthetic}).

\paragraph{On synthetic data.} \Cref{tab:bias-method} compares the three variants on a
\qty{100}{\hertz} grid. With noise but no lags, the IMU variant returns $\mtot$ within a few
percent at every window. The differentiated speed (B) underestimates it strongly at short
windows. That is the errors-in-variables effect of \cref{sec:sg-mechanics}: noise in the $\dot u$
regressor attenuates its coefficient. With the default lags, every variant misses $\mtot$ by
tens of percent, from \qty{-32}{\percent} (B, shortest window) to \qty{+102}{\percent} (A,
\qty{1}{\second}). The IMU does not remove the actuation lag between the logged PWM
and the real thrust. Because its acceleration edge is sharp, it even shows that lag more
clearly than the differentiated speed.

\begin{table}[htbp]
\centering
\caption{Relative error in percent of $\mtot$ and $X_u$ on a \qty{100}{\hertz} grid for the three
variants, three synthetic bags and three smoothing windows.}
\label{tab:bias-method}
\small
\begin{tabular}{@{}lc cc cc cc@{}}
\toprule
& & \multicolumn{2}{c}{A, current} & \multicolumn{2}{c}{B, consistent} & \multicolumn{2}{c}{C, IMU} \\
\cmidrule(lr){3-4}\cmidrule(lr){5-6}\cmidrule(l){7-8}
Bag & $t_\mathrm{w}$ (\unit{\second}) & $\mtot$ & $X_u$ & $\mtot$ & $X_u$ & $\mtot$ & $X_u$ \\
\midrule
\inputrows data/bias_method_rows.tex
\bottomrule
\end{tabular}
\end{table}
\FloatBarrier

\paragraph{On POLARIS.} The synthetic IMU is idealised: perfect timestamps, a constant and
exactly known pitch, no vibration and no lever arm. The real \code{/imu/data} of
\file{zermatt_rectangle_08} differs in three ways that matter.
\begin{enumerate}
  \item \textbf{Timestamps.} About \qty{86}{\hertz} on average, but delivered in bursts. The
        median spacing of the header stamps is \qty{0.1}{\milli\second}, with gaps of up to
        \qty{72}{\milli\second}. The stamps are batch arrival times, not sample times. Their
        error is of the same order as the timing errors that bias $\mtot$
        (\cref{tab:bias-grid,tab:bias-servo}).
  \item \textbf{Gravity leakage.} A pitch error $\delta\theta$ adds $g\sin\delta\theta$, which is
        \qty{0.17}{\meter\per\second\squared} per degree. A constant pitch error is absorbed by an
        offset term. However, \code{MOT_1} sits above the
        centre of mass (\cref{ch:symmetry}), so the pitch changes \emph{with the thrust}, and any
        error in that change leaks into $\dot u$ correlated with $T$. \Cref{tab:bias-leak}
        shows the effect.
  \item \textbf{Vibration.} About \qty{0.42}{\meter\per\second\squared} rms above
        \qty{5}{\hertz}. This is removed by any window of a few tenths of a second, so it is the
        least serious of the three.
\end{enumerate}

\begin{table}[htbp]
\centering
\caption{Relative error in percent of the IMU variant C on the ideal synthetic bag
(\qty{100}{\hertz} grid, \qty{0.5}{\second} window) when a pitch error proportional to the thrust,
reaching $\delta\theta_\mathrm{max}$ at the largest thrust, leaks gravity into the accelerometer.}
\label{tab:bias-leak}
\small
\begin{tabular}{@{}cc ccc@{}}
\toprule
$\delta\theta_\mathrm{max}$ (\unit{\degree}) & $g\sin\delta\theta_\mathrm{max}$ (\unit{\meter\per\second\squared}) & $\mtot$ & $X_u$ & $\Xuu$ \\
\midrule
\inputrows data/bias_leak_rows.tex
\bottomrule
\end{tabular}
\end{table}
\FloatBarrier

\begin{keyidea}
The IMU is the better $\dot u$ \emph{in principle}: nothing is differentiated, so the
noise-versus-window dilemma disappears. On POLARIS today it is not yet the better choice.
Its timestamps are unreliable at the tens-of-milliseconds level, and its gravity correction
depends on a pitch that changes with the thrust. Its best use now is to measure the actuation
delay from the sharp acceleration edge at a step (\cref{ch:timing}), and later to cross-check
$\mtot$ once the timestamps are fixed. Consistent filtering of the EKF speed remains the main
route, and output error (\cref{sec:output-error}) needs no $\dot u$ at all.
\end{keyidea}

\subsection{The servo log rate}

\code{servo_output_raw} is not the PWM itself but a sampled view of it. ArduSub changes the PWM
in its own fast loop as soon as the command changes. The topic only reports the current value
at its publishing rate $f_\mathrm{log}$. A step edge therefore appears in the log between $0$ and
$1/f_\mathrm{log}$ late, on average $1/(2f_\mathrm{log})$. This is a random lag of exactly the
kind that biases $\mtot$. The existing bags log it at \qty{2}{\hertz} (rectangle runs),
\qty{10}{\hertz} (June surface run) and \qty{20}{\hertz} (grid runs).

The synthetic bench cannot show this effect, because there the thrust follows the
\emph{logged} PWM (\cref{sec:synth-limits}). \file{gen_biasstudy.py} therefore uses a small
standalone simulation instead. The thrust switches at the command time, the PWM is logged at
$f_\mathrm{log}$ with a random phase, and variant B is fitted on a \qty{100}{\hertz} grid with a
\qty{0.5}{\second} window. Each rate is repeated with twelve random phases.

\begin{table}[htbp]
\centering
\caption{Relative error in percent of the fitted coefficients against the publishing rate
$f_\mathrm{log}$ of \code{servo_output_raw}, on an otherwise ideal surge run (mean over twelve random
log phases, and the range of $\mtot$).}
\label{tab:bias-servo}
\small
\begin{tabular}{@{}l c c c c@{}}
\toprule
$f_\mathrm{log}$ (\unit{\hertz}) & $\mtot$ (mean) & $\mtot$ (range) & $X_u$ & $\Xuu$ \\
\midrule
\inputrows data/bias_servo_rows.tex
\bottomrule
\end{tabular}
\end{table}
\FloatBarrier

\Cref{tab:bias-servo} shows the following.
\begin{itemize}
  \item At \qty{2}{\hertz}, $\mtot$ cannot be identified: the edge uncertainty of up to
        \qty{0.5}{\second} exceeds $\taueff$. The damping coefficients survive, because they come
        from the steady states, which do not depend on the edge timing.
  \item \qty{50}{\hertz} reduces the bias of $\mtot$ to about \qty{2}{\percent}, with a spread of a
        few percent. That is below the other error sources, so raising the rate to
        \qty{50}{\hertz} is worthwhile. \qty{100}{\hertz} gains little more.
  \item The numbers assume the dummy $\taueff \approx \qty{0.2}{\second}$. The bias scales
        roughly with $1/(f_\mathrm{log}\,\taueff)$, so a slower vehicle is less sensitive.
\end{itemize}
For the scripted manoeuvres there is also a way around the rate. The command scripts publish
\code{manual_control} at \qty{20}{\hertz} and mark every edge with \code{/sysid/event}, so the
command time of each edge is known to within a few milliseconds. With the command-to-PWM mapping
(open item 12 in \cref{ch:status}), the PWM can be reconstructed from the commands with exact
edges. ArduSub's internal latency then becomes part of the actuation delay, which is measured
anyway.

\subsection{Consequences}
\begin{enumerate}
  \item Filter every term of the surge equation with the same filter, and form $|u|u$ before
        filtering \eqref{eq:consistent}. This removes the window bias on clean data.
        \emph{Done in \file{fit_surge.py} on 16 October 2026}, see below.
  \item Use an analysis grid of at least \qty{100}{\hertz} (the cost is negligible).
        \emph{Done}: \code{RATE_HZ = 100}.
  \item Raise the \code{servo_output_raw} rate to \qty{50}{\hertz}, or reconstruct the PWM edges
        from \code{manual_control} and \code{/sysid/event}. Bags at \qty{2}{\hertz} are only
        good for the damping coefficients.
  \item The dominant bias remains timing, namely actuation lag and EKF lag, at
        \qtyrange{20}{100}{\percent} in $\mtot$ (\cref{tab:bias-method}). Measure the delay and
        compensate it (\cref{ch:synthetic,ch:timing}) before trusting any $\mtot$.
  \item Use the IMU to measure the delay now, and as a $\dot u$ source only once its timestamps
        are reliable.
\end{enumerate}

\begin{revised}[title={Applied to \file{fit_surge.py}, 16 October 2026}]
On request, items 1 and 2 were applied to \file{fit_surge.py}. The speed is smoothed once, and
$\dot u$ is \code{np.gradient} of the smoothed speed. $|u|u$ is formed from the raw speed and
then smoothed, the thrust passes through the same filter, and \code{RATE_HZ} went from 20 to
\qty{100}{\hertz}. The fit log records \code{filtering: consistent}. On the synthetic bags of
\cref{tab:bias-method} with a \qty{1}{\second} window, the script now returns the following.
\begin{center}
\small
\begin{tabular}{@{}lccc@{}}
\toprule
Bag & $\mtot$ & $X_u$ & $\Xuu$ \\
\midrule
ideal        & \qty{-1.6}{\percent} & \qty{-0.1}{\percent} & \qty{+0.1}{\percent} \\
noise only   & \qty{-5.6}{\percent} & \qty{+4.0}{\percent} & \qty{-0.9}{\percent} \\
all defaults & \qty{+60}{\percent}  & \qty{+15}{\percent}  & \qty{-7}{\percent} \\
\bottomrule
\end{tabular}
\end{center}
This matches variant B in \cref{tab:bias-method}. The remaining error on the ideal bag is the
grid quantisation at \qty{100}{\hertz}. The large error with all defaults is the timing (item 4),
which no filtering can remove.

Fixing only the quadratic term, while keeping the Savitzky--Golay derivative filter for $\dot u$,
is not enough. On the ideal bag with a \qty{1}{\second} window, $\mtot$ is then still
\qty{13}{\percent} too high at \qty{100}{\hertz} ($X_u$ \qty{+6}{\percent}). At
\qty{20}{\hertz} it is \qty{+6}{\percent}, but only because the grid bias of the opposite sign
hides part of the error. The mismatched derivative filter is therefore the larger of the two
inconsistencies.
\end{revised}

\section{Writing results back}

\code{model_io.update_parameters(\{"damping_linear.X_u": (12.3, 0.8)\}, bags=[...])} writes
value and uncertainty, sets \code{source: fitted}, bumps the patch version, stamps the
date, and appends the bag names to \code{fitted_from}. \code{Model.version_tag()} combines
the version, the sha256 of the model file and of every thrust table, and the git commit
(with \code{+dirty} if the file has uncommitted changes). Every propagation output and summary carries
this tag, so it can always be traced back to the model that produced it.

\begin{caution}
\code{update_parameters()} rewrites \file{model.yaml} with PyYAML, which \textbf{drops all comments}.
The comments in the \code{thrust:} and \code{allocation:} sections explain non-obvious
decisions, such as the frame conversion and the channel mapping. They should be preserved by
committing before the first automated update or by switching to a comment-preserving YAML library.
\end{caution}

\section{The fit log}
\label{sec:fit-log}

\begin{revised}[title={New section, added 16 October 2026}]
\file{results_log.py} and \file{results/fit_log.yaml} are new. \code{update_parameters()} gained
the argument \code{fit_id}, and \file{model.yaml} the field \code{model.fit_ids}.
\end{revised}

\file{model.yaml} only holds the \emph{current} values. When a coefficient is refitted, the
earlier value, the bag it came from and the settings that produced it are only in git
history, if at all. Fits that were never written into the model, such as trial runs,
rejected windows or a different smoothing window, leave no trace. The fit log closes that gap.
Every run of a fit script appends one entry to \file{scripts/sysid/results/fit_log.yaml}, which
is committed together with the code. The bags themselves are git-ignored, so each entry
identifies a bag by its path \emph{and} by a fingerprint.

\begin{table}[htbp]
\centering
\caption{Fields of one fit-log entry, as written by \code{results_log.log_fit()}.}
\label{tab:fit-log}
\small
\begin{tabularx}{\textwidth}{@{}l X@{}}
\toprule
Field & Content \\
\midrule
\code{id} & \code{fit-0001}, \code{fit-0002}, \dots, assigned in order \\
\code{date} & Local time with UTC offset \\
\code{script} & Repo-relative path of the fit script and its git state (\code{+dirty} if it has uncommitted changes) \\
\code{manoeuvre} & Free text, for example ``surge steps'' or ``coast-down'' \\
\code{bags} & Per bag: path, fit window \code{window_s} in seconds after the first EKF message
  (the \texttt{-{}-t-start}/\texttt{-{}-t-end} convention), size, and a fingerprint (sha256 over the
  size and the first and last \qty{1}{\mebi\byte}) \\
\code{model} & Path and \code{version_tag()} of the model file used, which includes the hashes
  of the thrust tables. A T200 fit and a T500 fit can therefore never be confused. \\
\code{settings} & Free dictionary: smoothing window, grid rate, \code{servo_delay_s}, \dots \\
\code{params} & Per parameter: value, uncertainty, unit. Names are free (for example \code{m_tot},
  which is not a \file{model.yaml} key). For model keys the unit is read from the model. \\
\code{metrics} & Free dictionary: $R^2$, residual standard deviation, sample count, condition number, \dots \\
\code{status} & \code{candidate} when written; may be edited by hand to \code{accepted} or \code{rejected} \\
\code{notes} & Free text \\
\bottomrule
\end{tabularx}
\end{table}
\FloatBarrier

\paragraph{Append only.} The code never rewrites the log. It reads the file only to count
the entries for the next id, and then appends one YAML list item. Comments and hand edits,
such as setting \code{status} to \code{rejected} with a comment giving the reason, therefore
survive. This is the opposite
of \code{update_parameters()}, which rewrites \file{model.yaml}.

\paragraph{Link to the model.} \code{update_parameters(..., fit_id="fit-0007")} appends the
id to \code{model.fit_ids}. A value in \file{model.yaml} can thus be traced to the entry
that produced it, and from there to the bag, the window and the settings.

\begin{incode}
At the end of a fit script, after the solve:
\begin{lstlisting}
from results_log import log_fit
fit_id = log_fit(__file__, [(args.bag, args.t_start, args.t_end)],
                 {"m_tot": (theta[0], se[0]),
                  "X_u": (theta[1], se[1]),
                  "X_uu": (theta[2], se[2])},
                 units={"m_tot": "kg", "X_u": "N s/m", "X_uu": "N s^2/m^2"},
                 manoeuvre="surge steps", model=args.model,
                 settings={"smooth_s": args.smooth_s, "rate_hz": args.rate},
                 metrics={"r2": r2, "n_samples": y.size})
\end{lstlisting}
The overview is printed by \code{python scripts/sysid/results_log.py}, which accepts
\texttt{-{}-param X\_u} (history of one parameter), \texttt{-{}-bag rectangle\_08} and
\texttt{-{}-show fit-0003}. The same \code{_append()} and \code{read_log()} are intended to carry
a validation log later.
\end{incode}

\begin{questions}
\begin{enumerate}
  \item Show that the least-squares residual is orthogonal to every column of $\vect{\Phi}$.
  \item Why would applying a causal (one-sided) low-pass filter to $u$ but not to $T$ bias $\mtot$?
  \item Name two advantages of output error over equation error for the step-from-rest manoeuvre.
\end{enumerate}
\end{questions}
